% TOP_PROBLEM: {"id":30007741,"problem_number":"LOCAL-30007741","title":"Healthy aging and demographic burdens","category_id":21,"category":{"id":21,"name":"economics","display_name":"Economics","description":"Open economic theory statements, published research agendas, and proposed scoped specifications; question provenance and historical priority are qualified per record.","slug":"economics","order_index":21,"created_at":"2026-10-08T00:00:00Z"},"status":"provenance_pending","external_url":null,"legacy_ids":[],"published":false,"statement":"Project the twenty-year fiscal effect of a specified improvement in functional health, separating healthier aging from survival changes and retirement responses.","economics":{"original_id":230,"field":"Education, families and demography","kind":"empirical","formalization_basis":"proposed_specification","source_status":"not_verified","priority_status":"not_established","priority_note":"No readable primary passage explicitly stating this precise open question was verified. A supporting paper, abstract, or topic citation is insufficient to establish origin or unresolved status.","earliest_verified_reference":null,"current_reference":{"authors":["Rainer Kotschy","David E. Bloom","Andrew J. Scott"],"title":"On the Limits of Chronological Age","year":2024,"url":"https://docs.iza.org/dp17427.pdf","doi":null,"locator":"Introduction and empirical results; IZA Discussion Paper 17427","evidence_summary":"The paper documents limits of chronological-age proxies; it does not itself establish the broader fiscal-offset question as an open problem."},"formalization_note":"Proposed scoped research specification. The original catalogue prompt was a synthesis, and no canonical conjecture or verified original formulation is claimed. The supporting source, when retained, establishes context or existing results rather than this exact open question."},"statusQualification":"Provenance pending: an explicit published statement of this exact open question was not verified. This is a catalogue-authored candidate specification, not a certified open conjecture. Proposed scoped research specification. The original catalogue prompt was a synthesis, and no canonical conjecture or verified original formulation is claimed. The supporting source, when retained, establishes context or existing results rather than this exact open question.","statusReviewedAt":"2026-10-08"}
% FIRST_POSTED: 2026-10-08
% ENTRY_KIND: statement_only
% !TeX program = lualatex
\documentclass[11pt]{article}
\usepackage{fontspec}
\usepackage[a4paper,margin=25mm]{geometry}
\usepackage{amsmath,amssymb,mathtools}
\usepackage{xurl}
\usepackage[hidelinks]{hyperref}
\newcommand{\sref}[2]{\hyperlink{source:#1:#2}{[S#2]}}
\setlength{\emergencystretch}{3em}
\begin{document}
% problemId:30007741
\section*{Healthy aging and demographic burdens}\label{30007741}
\subsection{Definitions and mathematical statement}
Consider residents aged fifty through eighty-five in one high-income country over a twenty-year fiscal projection. Let \(a=1\) denote a specified improvement in age-specific functional health relative to baseline \(a=0\), preserving a stated mortality path or explicitly modeling its change. For individual \(i\), define annual labor earnings \(Y_{it}(a)\), public health and care expenditure \(H_{it}(a)\), pension transfers \(P_{it}(a)\), and taxes \(T_{it}(a)\). For discount factor \(\beta\in(0,1)\), define net fiscal burden
\[B(a)=E\!\left[\sum_i\sum_{t=0}^{20}\beta^t\{H_{it}(a)+P_{it}(a)-T_{it}(a)\}\right].\]
Determine \(B(1)-B(0)\) under stated retirement rules, wage responses, caregiving arrangements, and demographic projections. Estimate health’s effects on work and care needs rather than infer behavior from chronological age alone. Separate healthier aging from longer survival, and report distributional household welfare alongside the fiscal comparison. Causal identification must address selection in health and employment; observational age profiles are insufficient. The cited paper demonstrates limits of chronological-age proxies, but does not explicitly establish this precise fiscal-offset question as open. This is therefore a proposed empirical projection with transparent behavioral assumptions and no verified original priority.

\par\textbf{Formulation and scope.} Proposed scoped research specification. The original catalogue prompt was a synthesis, and no canonical conjecture or verified original formulation is claimed. The supporting source, when retained, establishes context or existing results rather than this exact open question.

\par\textbf{Open-question provenance.} An explicit published open-question statement was not verified. Sources below are supporting literature and must not be treated as a first listing.

\par\textbf{Historical priority.} No readable primary passage explicitly stating this precise open question was verified. A supporting paper, abstract, or topic citation is insufficient to establish origin or unresolved status.

\subsection{Short English statement}
Project the twenty-year fiscal effect of a specified improvement in functional health, separating healthier aging from survival changes and retirement responses.

\subsection{Sources}
\begin{itemize}\raggedright
\item[\textnormal{[S1]}]\hypertarget{source:30007741:1}{} Rainer Kotschy, David E. Bloom, Andrew J. Scott. 2024. On the Limits of Chronological Age. Introduction and empirical results; IZA Discussion Paper 17427.
\url{https://docs.iza.org/dp17427.pdf}.
{\small\textit{Supporting or current literature; see evidence qualification. The paper documents limits of chronological-age proxies; it does not itself establish the broader fiscal-offset question as an open problem.}}
\end{itemize}
\end{document}
